% ===== §8 Aesthetic Education =====
\section{Aesthetic Education: Cultivating the Judgement of the Good}
\label{sec:aesthetic-education}

\noindent
If the construction of the field is an aesthetic problem (\xref{sub:constr-thesis}), and the perception of what arises within it is an aesthetic skill (\xref{sec:aesthetic-perception}), then the cultivation of the capacity for both is an \emph{aesthetic education}. This is the philosophical summit of the paper, and also its most exposed wager, for it will argue that aesthetic education is not a refinement added to ethical life but, at the limit, the most fundamental form of ethical education there is. The section advances three theses of increasing radicality. The first, already stated, is that the construction of the field is aesthetic. The second distinguishes two kinds of beauty, the received and the enacted, and locates the second at the root of the long-noted kinship between practical ethics and the aesthetic. The third, the most radical, is that ethical judgement, at the point where rules give out, operates aesthetically, so that the good cannot finally be separated from the beautiful, and the deepest ethical education is aesthetic. The section defends the third thesis in a guarded form, confronts it with its strongest objectors, surveys the classical and modern theories of aesthetic education, descends to the question of how the relevant capacities are actually cultivated between two people, and closes on the difficulty that the whole edifice rests upon: that beautiful things are difficult.

% ============================================================
\subsection{Why aesthetic education, and not simply ethical education}
\label{sub:ae-why}

The claim that the cultivation of the good is an aesthetic rather than a narrowly ethical education depends on a result established in the previous paper and reactivated here: that the good is not programmable \citep{paperxx}. If the good of an intimate relation could be reduced to a set of rules, do this, refrain from that, follow the procedure, then its cultivation would be a matter of teaching the rules, and ethical education would be the transmission of a code. But the good of a relation cannot be so reduced; it requires, in each concrete case, a judgement of what is fitting here, now, with this person, that no rule determines, because the concrete case always exceeds the rule's generality. This was the lesson of the previous paper's closing recognition that the good must be judged afresh in the flesh of each particular relation, and it has a direct consequence for education: if the good is not a code, its cultivation cannot be the transmission of a code, and must instead be the cultivation of the \emph{judgement} that operates where the code gives out.

What is the structure of that judgement? It is the judgement of what is fitting in a concrete case that no rule subsumes, a holistic, situated grasp of the right proportion, the apt response, the fitting measure, exercised on a particular that no universal exhausts. And this, the section will argue, is the structure of aesthetic judgement. The judgement that catches the right measure in a concrete case, where no rule can supply it, is the same judgement in the construction of a field, the perception of a contingency, and the response to an ethical situation; it is, in each, a felt grasp of the fitting that is irreducible to procedure. If that is so, then the cultivation of the good is the cultivation of this judgement, and the cultivation of this judgement is aesthetic education, the education of the capacity to perceive and respond to the fitting where no rule can be followed. This is why the education the good requires is aesthetic and not merely ethical in the narrow, rule-transmitting sense: because the good itself, at its un-programmable core, requires a judgement whose structure is aesthetic.

% ============================================================
\subsection{Thesis Two: receptive beauty and enactive beauty}
\label{sub:ae-thesis2}

Before the most radical thesis, a distinction must be drawn that the discussion has so far left implicit, and that proves to be the hinge between the aesthetic and the ethical. Beauty, as it figures in this paper, is of two kinds.

\begin{thesis}{Thesis Two.}
Beauty divides into the \emph{receptive} and the \emph{enactive}. Receptive beauty is beauty as received, the beauty of the wildflower, the landscape, the work of art, met in the aesthetic perception that beholds it. Enactive beauty is beauty as done, the beauty of an action, a response, a way of treating another, made beautiful in the doing. The first is perceived; the second is enacted. And it is enactive beauty that grounds the long-observed kinship between practical ethics and the aesthetic, for to act so that one's treatment of another is itself beautiful is to enact the good as beauty.
\end{thesis}

\noindent
Receptive beauty is the beauty the previous section was concerned with: the contingent beauty caught in perception, the flower seen and the landscape beheld, beauty as something that arrives and is received. It is the object of the aesthetic perception of contingency, and its cultivation is the cultivation of the receptive capacities, the trained eye, the awakened \emph{xing}, the openness to what the world offers. But there is a second beauty that is not received but \emph{made}: the beauty of a gesture, a response, a way of meeting another person, that is beautiful in the way it is done. A response to a partner's distress can be beautiful, apt, proportioned, generous in exactly the right measure, neither too much nor too little, fitting the moment as a perfect line fits a poem. A way of treating another can have grace, the way a movement or a phrase can have grace. This is enactive beauty: not beauty beheld but beauty enacted, the beauty of the deed.

The distinction matters because enactive beauty is where the aesthetic and the ethical meet, and meet not by analogy but in substance. To act well toward another, to respond to their need with exactly the right measure, to meet their vulnerability with exactly the fitting grace, to treat them with the aptness that the situation calls for and no rule prescribes, is at once an ethical achievement and an aesthetic one, and the two are one rather than two descriptions of the achievement. The rightness of the response is its fittingness, and its fittingness is its beauty; the well-met moment is the beautifully-met moment, and there is no separating, in it, the ethical aptness from the aesthetic grace. This is why practical ethics has so often been felt to have an aesthetic dimension, why we speak of a response as graceful, of a life as beautiful, of conduct as fitting; the kinship is a recognition that enacted goodness is enacted beauty rather than a metaphor borrowed from art, that to do the fitting thing is to make a small beautiful thing of one's action. Enactive beauty is the practical beauty of the well-done deed, and its cultivation is the cultivation of the capacity to act beautifully, which is, the next thesis will argue, indistinguishable at its root from the capacity to act well.

% ============================================================
\subsection{Thesis Three: ethical judgement is, at its limit, aesthetic}
\label{sub:ae-thesis3}

The third thesis is the paper's most radical, and it is advanced with the care its radicality demands. It is stated first in a guarded form, which the paper defends; then in an extreme form, which the paper entertains as a genuine temptation and holds open rather than asserts.

\begin{thesis}{Thesis Three (guarded form).}
At the point where rules are exhausted by the concrete, which is every point at which a genuinely particular ethical situation must be met, ethical judgement operates aesthetically. Practical wisdom, at its summit, coincides with aesthetic judgement: both are the holistic, non-conceptual grasp of what is fitting in a concrete case that no rule subsumes. The cultivation of the good is therefore, at its core, the cultivation of aesthetic judgement, and the deepest ethical education is an aesthetic education.
\end{thesis}

\noindent
The argument for the guarded form proceeds in two steps, the first structural and the second concerning purity. The structural step is this. Ethical life cannot be exhaustively governed by rules, because rules are general and situations are particular, and the particular always in principle exceeds the general; there is always a gap between the rule and the case, which judgement must cross. What crosses it? Not another rule, for that would only reproduce the gap; but a capacity to grasp, in the concrete situation, what is fitting, a capacity Aristotle named \emph{phronesis}, practical wisdom, the perception of the right thing to do in the particular case, which he was careful to say could not be reduced to rules and belonged to perception rather than to demonstration \citep{aristotle2000}. Now consider the structure of aesthetic judgement as Kant analysed it: the judgement of beauty is, for Kant, a \emph{reflective} judgement, one that finds the universal in the particular rather than subsuming the particular under a given universal; it is non-conceptual, it does not apply a determinate concept, yet it claims a kind of universality; it is the free play of the faculties grasping a purposiveness without a determinate purpose, a fittingness with no rule that the fittingness instances \citep{kant1987}. The two structures, \emph{phronesis} and aesthetic judgement, are, the thesis claims, the same structure: both are the non-conceptual grasp of a fitting particular that no rule subsumes, the finding of the right measure in the concrete where no determinate universal can be applied. \emph{Phronesis} at its summit and aesthetic judgement coincide: the perception of the fitting in the ethical case and the perception of the fitting in the aesthetic case are one capacity, exercised on different materials. This is the bridge Kant's third \emph{Critique} was written to build, the judgement that mediates between the realm of nature and the realm of freedom \citep{kant1987}, and the thesis takes it at its word: the faculty that judges beauty is the faculty that, at its limit, judges the good.

The second step concerns purity, and it is the heart of the matter and the author's own deepest reason for the thesis. Consider the difference between two ways of doing the right thing. In the first, one does the right thing because a rule, an obligation, an internalised norm, requires it: one has been formed to act so, and one acts so because one is so formed. In the second, one does the right thing because one perceives, freshly and for oneself, that it is the fitting thing to do, because one sees its fittingness, is moved by it, and enacts it out of that perception, as one might be moved to complete a line of verse in the only way that fits. The two may issue in the same outward act; but they differ in their source, and the difference is the difference between a structured second nature and a free response. The rule-governed good is \emph{structured humanity}: a humanity formed by the internalisation of external norms, acting rightly because it has been built to, its rightness the operation of an installed structure rather than a free perception. Such goodness is reliable, and it is not nothing, a world of people reliably formed to act well is far better than its opposite, but it is not \emph{pure}, for it is heteronomous in the deepest sense: the source of the action is the installed structure, the internalised rule, the formed disposition, and not the free perceiving subject. The aesthetically-grounded good is different. To act rightly because one perceives, oneself and freshly, the fittingness of the act, to be moved by the beauty of the fitting response and to enact it out of that being-moved, is to act from a source that is one's own free perception rather than an installed structure; and only such action, the thesis holds, has purely ethical character, because only in it is the subject the genuine origin of the good it does, rather than the executor of a structure built into it. The rule-governed good is safe, reliable, and structured; the aesthetically-grounded good is dangerous, unreliable, and free, and only the free is pure. This is why, on the deepest version of the thesis, the truly ethical can only be the aesthetic: because only a goodness that springs from the free perception of the fitting, rather than from the operation of an internalised rule, is the goodness of a free subject rather than the functioning of a structured one.

\begin{thesis}{Thesis Three (extreme form, entertained not asserted).}
Every ethical question is, at bottom, an aesthetic question. There is, finally, no good that is not the beautiful enacted, and no ethical judgement that is not, at its root, a perception of the fitting, so that ethics is a province of aesthetics, and the good a name for the beautiful in the realm of action.
\end{thesis}

\noindent
The extreme form is a genuine temptation, and the paper neither asserts nor dismisses it. It is the thought that the whole distinction between ethics and aesthetics is, at the limit, unsustainable, that what we call the good just is the beautiful as it appears in conduct, and that the apparent difference between them is an artefact of treating ethics as a matter of rules and aesthetics as a matter of taste, a difference that dissolves once one sees that ethics at its core is also a matter of the perception of the fitting. The paper marks this as the radical horizon of its argument and declines to assert it, for reasons the next subsection makes clear: the extreme form, unguarded, courts a danger severe enough that responsibility requires holding it at the level of a temptation rather than a thesis. But it records it, because the guarded form points toward it, and an honest account does not conceal the precipice to which its own argument leads.

% ============================================================
\subsection{The strongest objections, and the reply}
\label{sub:ae-objections}

A thesis this radical must face its strongest objectors at full strength, and there are two whose force is sufficient to require, not refutation, but the guarding that the thesis has already built into itself. The series' practice has been to let the hardest objection stand at full force \citep{paperxx}; the practice is kept here.

\paragraph{The Kantian objection: the heteronomy of inclination.}
The first objection comes from Kant's own ethics, and it is sharp precisely because the thesis leaned on Kant's aesthetics. For Kant the moral worth of an action lies in its being done from duty, from respect for the moral law, and \emph{not} from inclination, however amiable; an action done because one is moved to it, because one feels its fittingness or is drawn by its beauty, is for Kant exactly \emph{not} a moral action in the full sense, but at best a fortunate \emph{pathological} one, done from feeling rather than from the rational recognition of duty. On this view the thesis has it exactly backwards: to ground the good in the aesthetic perception of the fitting, in being-moved by the beauty of the response, is to ground it in inclination, and inclination is precisely what Kant excludes from the purely moral. The aesthetically-grounded good, far from being the pure good, would be on Kantian grounds the impure one, the good done from feeling rather than from law.

The reply does not deny the Kantian point but relocates the disagreement. The objection assumes that the alternative to acting-from-feeling is acting-from-the-rational-law, and that purity lies with the latter. But the thesis's purity argument was directed at exactly this: the rational law, internalised and operative as the source of one's action, is itself a structure, the structure of duty installed in the subject, and action from it is, in the sense the thesis specified, the operation of an installed structure rather than the free response of the perceiving subject. The thesis and Kant agree that inclination as mere pathological impulse is no ground for the good; they disagree about whether the free perception of the fitting is ``inclination'' in that sense. The thesis holds that it is not, that the perception of the fitting, the being-moved by the beauty of the right response, is a \emph{judgement} rather than a blind pathological pull, the reflective grasp of a fittingness, which has more in common with Kant's own aesthetic judgement (universal, disinterested, reflective) than with the inclinations he excluded. The disagreement, then, is whether the faculty of reflective judgement, which Kant himself placed between nature and freedom, can be the source of moral action; Kant says the source must be pure practical reason, the thesis says the source can be the reflective perception of the fitting. This is a real and unresolved disagreement, and the paper does not pretend to settle it against Kant; it claims only that the thesis is not refuted by the charge of grounding the good in mere inclination, because the perception of the fitting is a judgement rather than mere inclination, and that whether such judgement can be a pure ground of the good is exactly the open question the thesis raises against the Kantian orthodoxy, not a confusion the orthodoxy disposes of.

\paragraph{The objection from the aestheticisation of ethics: the fascist spectre.}
The second objection is graver, because it is historical and not merely conceptual, and because the previous paper already conceded its force in another form \citep{paperxx}. To make the good a matter of the beautiful, to aestheticise ethics, is to invite the most dangerous confusion of the last century: the aestheticisation of politics that Benjamin diagnosed at the heart of fascism, the treatment of the political as a work of art to be composed, the people as material to be shaped into a beautiful whole, the leader as artist and the nation as artwork \citep{benjamin1969}. If the good is the beautiful enacted, what prevents the beautiful-but-monstrous, the magnificently staged rally, the aesthetically perfect and morally abominable order, the cruelty committed with style? The previous paper conceded exactly this point in its own register, granting that a relation could be formally beautiful, could have the geometry of the good cycle, and yet be unjust, so that aesthetic form alone does not secure goodness \citep{paperxx}. The objection here is the same at a higher pitch: the aestheticisation of ethics threatens to license the beautiful atrocity, to make style a sufficient credential for the good, and history shows where that leads.

This is the objection that forces the thesis into its guarded form and keeps the extreme form a temptation rather than an assertion. The reply has three parts, and none of them fully dissolves the danger, which is as it should be. First, the thesis distinguishes receptive and enactive beauty, and the fascist aestheticisation is an aestheticisation of the \emph{spectacle}, of receptive beauty, the beauty beheld by the masses in the staged rally, not of the enactive beauty of how persons are treated; and enactive beauty, the beauty of the fitting treatment of \emph{this} person, is precisely what the beautiful atrocity lacks, since the atrocity treats persons as material for a spectacle rather than meeting them in the aptness their particularity demands. The fascist spectacle is receptively magnificent and enactively monstrous; the distinction the thesis draws is exactly the one that the beautiful atrocity violates. Second, the enactive beauty the thesis prizes is inseparable, by the analysis of the whole series, from the non-extractive, non-dominative treatment of the other, the fitting response is the one that meets the other as an end and returns rather than retains \citep{paperxx,paperxvii}, so that a ``beautiful'' treatment that reduced the other to material would, on the thesis's own account of enactive beauty, not be beautiful but its counterfeit, the meretricious mistaken for the beautiful. Third, and most honestly, the danger is real and the guarding is necessary precisely because the extreme form, unguarded, could not make these distinctions stick: if every ethical question simply \emph{were} an aesthetic one, with no further criterion, then the magnificent and the monstrous could not be told apart by anything but taste, and taste can be trained to monstrosity. It is exactly to keep the criterion of enactive beauty, the fitting, non-extractive meeting of the particular other, from collapsing into mere spectacular magnificence that the thesis must be held in its guarded form, in which the aesthetic judgement that grounds the good is the judgement of \emph{enactive} fittingness and not of \emph{receptive} splendour. The fascist spectre is the standing proof that the extreme form is too dangerous to assert; it is the reason the paper entertains that form as a precipice rather than standing on it. \emph{The beautiful that grounds the good is the beauty of the fitting deed, not the beauty of the magnificent show}, and the whole weight of the thesis rests on a distinction that history shows to be perpetually under threat of erasure.

% ============================================================
\subsection{The classical theories of aesthetic education}
\label{sub:ae-classical}

The thought that the cultivation of the human being passes essentially through the beautiful is old and deep, and the thesis of this paper is its inheritor rather than its inventor. The classical theories are surveyed in Table~\ref{tab:ae-classical}; the text draws out the four that bear most directly on the argument.

The foundational modern text is Schiller's \emph{On the Aesthetic Education of Man}, and it is foundational precisely because it makes the claim this paper makes: that the way to a free and whole humanity runs through the aesthetic \citep{schiller1967}. Schiller's problem was the one this paper has met in another form, how a being divided between sensuous inclination and rational law might become free and whole rather than either the slave of impulse or the executor of an installed law. His answer was the \emph{play-drive} (\emph{Spieltrieb}), the faculty that arises in the free play between the sensuous and the rational and that is exercised paradigmatically in the experience of beauty; in aesthetic play the human being is, Schiller said, most fully human, because there alone the two halves of its nature are reconciled in freedom rather than one subordinated to the other. This is, in a different vocabulary, exactly the thesis's contrast between the structured humanity of the installed rule and the free humanity of the aesthetic response: Schiller's aesthetic education is the cultivation of the freedom that lies beyond both impulse and law, in the play where the human being is whole. The paper's purity argument is Schiller's insight in the register of ethics, that only the aesthetically free response is the response of a whole and free human being, rather than of one half of a divided one obeying the other.

The decisive analysis of the faculty is Kant's, already invoked: the \emph{Critique of Judgment} supplies the structure of the reflective aesthetic judgement on which the whole thesis turns, and it is Kant who frames that faculty as the bridge between nature and freedom, the very mediation the thesis exploits \citep{kant1987}. Aristotle supplies the ethical half of the bridge, the \emph{phronesis} whose structure the thesis identifies with aesthetic judgement \citep{aristotle2000}. And the Confucian tradition supplies what the Western sources lack: an explicit, institutionalised programme of aesthetic education as the core of ethical formation. The \emph{Analects} states it with a concision the West never matched: \zh{兴于诗，立于礼，成于乐}, one is aroused by poetry, established by ritual, and \emph{completed by music} \citep{confucius_analects}. Here the cultivation of the person begins in the affective stirring of poetry (the \emph{xing} of \xref{sub:ap-moments}) and \emph{culminates} in music, the ethical formation of the human being completed in an aesthetic accomplishment, the harmony of the well-formed character figured as, and achieved through, the harmony of music. This is the classical locus of the paper's two beauties: \emph{xing}, the arousal by poetry, is receptive beauty at the root of formation; \emph{cheng yu yue}, completion in music, is enactive beauty as the summit of the formed character, the life made harmonious as a piece of music is made harmonious. The Confucian sequence is the thesis of this section stated two and a half millennia early: ethical cultivation begins and ends in the aesthetic.

\begin{table}[ht]
\centering
\small
\begin{tabularx}{\linewidth}{L{2.7cm} L{3.5cm} Y}
\rowcolor{tablehead}\lrhead{Source} & \lrhead{Core concept} & \lrhead{Bearing on the thesis} \\
\midrule
Schiller \citep{schiller1967,schiller_naive} & The play-drive (\emph{Spieltrieb}); beauty as the reconciliation of sense and reason & The freedom beyond both impulse and law is reached through the aesthetic --- the paper's purity argument in its original form. \\
Kant \citep{kant1987} & Reflective aesthetic judgement; judgement as the bridge of nature and freedom & Supplies the structure of the faculty the thesis identifies with \emph{phronesis}; the mediation the thesis exploits. \\
Aristotle \citep{aristotle2000} & \emph{Phronesis}: perception of the fitting in the particular, irreducible to rule & The ethical half of the bridge; practical wisdom as non-conceptual grasp of the fitting. \\
Confucius \citep{confucius_analects} & \zh{兴于诗，立于礼，成于乐}: aroused by poetry, completed in music & Ethical formation begins (\emph{xing}, receptive) and culminates (music, enactive) in the aesthetic. \\
Plato \citep{plato1926} & The difficulty of the beautiful; beauty and the good ascending together & \el{χαλεπὰ τὰ καλά}: the beautiful resists definition --- the un-programmability at the section's close. \\
Daoist \citep{laozi1963} & \zh{虚}; the use of emptiness; the uncarved & The aesthetic education of the perception of spaciousness and the un-forced. \\
\end{tabularx}
\caption{Classical theories of aesthetic education and their bearing on the paper's three theses.}
\label{tab:ae-classical}
\end{table}

% ============================================================
\subsection{The modern theories of aesthetic education}
\label{sub:ae-modern}

The modern theories, surveyed in Table~\ref{tab:ae-modern}, develop the classical inheritance in three directions that matter to the paper: the democratisation of aesthetic experience, its critical and emancipatory edge, and its concrete pedagogy.

Dewey's \emph{Art as Experience} performs the decisive democratising move, and it is the modern source closest to the paper's own argument \citep{dewey1934}. Against the confinement of the aesthetic to museums and masterpieces, Dewey located aesthetic experience in the consummation of ordinary experience, in any stretch of lived experience that achieves a certain wholeness, completion, and felt unity, which he called ``an experience.'' The roadside flower, caught and consummated into a complete small experience, is aesthetic in exactly Dewey's sense; the aesthetic perception of contingency is the Deweyan finding of the aesthetic in the texture of ordinary life, and the paper's insistence that the good of the everyday is available without grand apparatus is Dewey's democratisation carried into the ethics of intimacy. Cai Yuanpei's programme of ``replacing religion with aesthetic education'' carries the democratising impulse into a whole national pedagogy, proposing aesthetic cultivation as the universal, secular path to the formation of a free and elevated humanity, the institutional dream of exactly the thesis this paper argues \citep{caiyuanpei2015}.

The critical edge is supplied by the Frankfurt tradition. Adorno's aesthetics locates in the genuine aesthetic experience a resistance to instrumental reason, the aesthetic as the comportment that refuses to reduce its object to a use, that beholds rather than exploits, and so stands against the total administration of the world by the logic of means and ends \citep{adorno1997}. This gives the paper's aesthetic perception its critical-political dimension, taken up in the next section: to perceive aesthetically, to behold the contingent for its own sake rather than for its use, is already a small refusal of the instrumentalisation that the administered world imposes. Rancière sharpens this into an explicit politics: the aesthetic is the ``distribution of the sensible,'' the prior partition of what can be perceived, by whom, as significant, so that to alter what is perceptible, to make the overlooked perceptible, is a political act, and aesthetic education, the cultivation of new capacities of perception, is the redistribution of the sensible \citep{ranciere2004}. The cultivation of the aesthetic perception of contingency is, in Rancière's terms, the acquisition of a capacity to perceive what the dominant distribution renders imperceptible, a politics of perception, whose stakes the following section will make concrete.

\begin{table}[ht]
\centering
\small
\begin{tabularx}{\linewidth}{L{2.7cm} L{3.5cm} Y}
\rowcolor{tablehead}\lrhead{Source} & \lrhead{Contribution} & \lrhead{Bearing on the paper} \\
\midrule
Dewey \citep{dewey1934} & ``An experience''; the aesthetic in the consummation of ordinary life & Democratises the aesthetic; the roadside flower is aesthetic in Dewey's exact sense. \\
Cai Yuanpei \citep{caiyuanpei2015} & Aesthetic education in place of religion; universal secular formation & The institutional dream of the thesis; aesthetic cultivation as the path to a free humanity. \\
Adorno \citep{adorno1997} & The aesthetic as resistance to instrumental reason & Beholding-not-using as a refusal of administration; the critical edge of perception. \\
Rancière \citep{ranciere2004} & The distribution of the sensible; aesthetics as politics & Cultivating perception redistributes the perceptible --- a politics of perception. \\
Murdoch \citep{murdoch1970} & Attention; the just and loving gaze & The moral work of attention is aesthetic in structure --- seeing truly is loving rightly. \\
Scarry \citep{scarry1999} & Beauty as prompting justice; the decentred, fair regard & Receptive beauty trains the fair attention that the just response requires. \\
Modern pedagogy \citep{montessori1967,steiner1996,greene2001,eisner2002} & Sensory education; artistic formation; aesthetic education as practice & Concrete methods for cultivating perception --- how the capacities are actually trained. \\
\end{tabularx}
\caption{Modern theories of aesthetic education and their bearing on the paper.}
\label{tab:ae-modern}
\end{table}

Two further modern sources deserve note for joining the aesthetic to the ethical in exactly the paper's way. Iris Murdoch made \emph{attention}, the just and loving gaze, the patient regard that sees the other truly, the centre of the moral life, and her attention is aesthetic in structure: it is the disinterested, careful beholding that the aesthetic perception of a thing requires, turned upon a person, so that to see the other justly is to behold them with the regard one gives a thing of beauty \citep{murdoch1970}. Elaine Scarry argued that the experience of beauty trains the fair, decentred attention that justice requires, that to be stopped and decentred by a beautiful thing is to practice the very unselfing that just regard for others demands \citep{scarry1999}. Both are, in the paper's terms, theorists of how receptive beauty schools enactive beauty: how the trained perception of the beautiful becomes the capacity for the fitting, attentive, just response. And the concrete pedagogies, Montessori's education of the senses, Steiner's artistic formation of the whole child, the aesthetic-education programmes of Greene and Eisner, supply what a paper on cultivation must not omit: actual methods by which the capacities of aesthetic perception are trained, evidence that the aesthetic education the thesis requires is a practiced art with a literature of its own, rather than a pious wish \citep{montessori1967,steiner1996,greene2001,eisner2002}.

% ============================================================
\subsection{How it is cultivated: the mutual aesthetic education of two}
\label{sub:ae-cultivation}

The theories describe the cultivation of a person; this paper's concern is the cultivation of a ``we,'' and the distinctive claim of this subsection is that between two people aesthetic education is \emph{mutual}, that each is the other's teacher, and the relation itself a school. This is not a metaphor. The capacities at issue, the perception of contingent beauty, the receptivity of \emph{xing}, the judgement of the enactively fitting, are cultivated, between two people, largely \emph{by each other}, in three concrete ways.

The first is the pedagogy of \emph{look}. When one partner catches a contingent beauty and turns it toward the other, \emph{look}, they do not merely share a moment; they \emph{teach a way of seeing}. The turning-toward is a demonstration: it shows the other what is worth catching, trains the other's eye toward the kind of thing one has learned to see, awakens in the other a \emph{xing} they might not have had. Over time, two people who habitually turn contingent beauties toward each other educate each other's perception, each becoming more able to catch what the other catches, until they come to share a trained sensibility, a common eye for the beautiful, built up from a thousand acts of mutual showing. The relation becomes a school of perception in which each is continually teaching the other to see, and the curriculum is the contingent beauty of the shared world.

The second is the pedagogy of shared retelling, which is also the reproduction of experience (\xref{sub:reproduction-retelling}) seen from the side of aesthetic education. To retell a shared journey, to return together to the beauty once caught, is to re-cultivate the perception of it, to deepen, in the retelling, the capacity to perceive what the experience held, and to refine together the sensibility that first caught it. Shared retelling is mutual aesthetic education exercised upon the past: the joint re-perceiving of what was lived, by which the two refine together their common capacity to perceive. This connects the paper's political economy to its aesthetics: the reproduction of experience is, at the same time, the mutual cultivation of the sensibility that perceives experience, so that the labour that conserves the relation's wealth is the same labour that educates its perception.

The third is the pedagogy of enacted beauty: each partner teaches the other the enactively fitting by enacting it. When one meets the other's vulnerability with exactly the right measure, responds to their distress with exactly the fitting grace, the other learns, not by precept but by undergoing, what the fitting response is, and is formed toward the capacity to give it in turn. To be treated beautifully is to be educated in the enactment of beauty; the well-met partner learns, from being well met, how to meet. The relation is thus a reciprocal school of enactive beauty, each forming the other's capacity for the fitting response by the gift of fitting responses received. And this is the deepest sense in which the relation is the mutual cultivation of the good: by treating each other beautifully, two people educate each other in the enactment of the good, so that the relation becomes, over a lifetime, a mutual aesthetic education in the most consequential of all the arts, the art of meeting another person fittingly, which is the same, by the argument of this section, as the art of treating them well.

% ============================================================
\subsection{\texorpdfstring{\el{χαλεπὰ τὰ καλά}}{Chalepa ta kala}: that the beautiful is difficult}
\label{sub:ae-difficult}

It remains to name the difficulty on which the whole edifice rests, and it is fitting to name it with the oldest words for it. In Plato's \emph{Hippias Major}, Socrates and the sophist Hippias attempt to say what the beautiful is, and fail; Hippias, who can speak fluently and at length on every other subject, is reduced, on this one, to a succession of definitions that Socrates dismantles one by one, until the dialogue ends not in an answer but in the proverb Socrates offers in place of one: \el{χαλεπὰ τὰ καλά}, beautiful things are difficult \citep{plato1926}. It is a remarkable ending. The man who could define anything could not define the beautiful; the dialogue that sought the essence of beauty found only its difficulty. And the difficulty is not accidental but essential: the beautiful resists definition because it is not the kind of thing a definition can capture, not a property specifiable by a rule, but exactly the un-ruled fittingness that the whole of this section has been concerned with. \el{χαλεπὰ τὰ καλά} is the most ancient name for the un-programmability of the beautiful, and by the thesis of this section, of the good.

This difficulty is its confirmation rather than a defect of the account, and it explains the section's central claim about education. Because the beautiful cannot be reduced to a rule, it cannot be \emph{taught}, not in the sense in which a rule is taught, by transmission of a content to be applied. It can only be \emph{cultivated}, developed, in oneself and in another, through long exposure, practice, attention, and the slow refinement of a perception, the way taste is cultivated rather than taught, the way an eye is trained rather than instructed. This is exactly why the cultivation of the good must be an aesthetic \emph{education} in the deep sense and not a moral \emph{instruction}: because the good, like the beautiful and for the same reason, cannot be transmitted as a rule but can only be cultivated as a sensibility. And it is why the cultivation is so difficult, so slow, so resistant to method, so dependent on the long mutual schooling of two people who teach each other to see and to meet over a whole shared life. \el{χαλεπὰ τὰ καλά}: the beautiful is difficult, the good is difficult, and they are difficult in the same way and for the same reason, because neither is a rule to be followed but a fitting to be perceived, and the perception of the fitting is the most difficult and the most uninstructable of all the human capacities, the one that can only be grown. That it is difficult is not a reason to doubt that it is the heart of the matter; it is the sign that one has reached the heart of the matter, the place where Socrates himself was reduced to a proverb, and where this paper, having followed the good to its root in the beautiful, can do no more than name the difficulty and turn to the cultivation that is the only response to it.
